% 第1章 固体中的电子相互作用
\chapter{固体中的电子相互作用}

\section{单电子理论}

在周期势中运动的单个电子由能带结构方程描述：

\begin{equation*}
H ^ {0} \phi_ {\mathbf{k} s} (\mathbf{x}) \equiv \left[ - \frac {\hbar^ {2}}{2 m} \nabla^ {2} + V ^ {\mathrm{ion}} (\mathbf{x}) \right] \phi_ {\mathbf{k} s} (\mathbf{x}) = \epsilon_ {\mathbf{k}} \phi_ {\mathbf{k} s} (\mathbf{x}),
\tag{1.1}
\end{equation*}

其中 $\phi_{\mathbf{k}s}$ 与 $\epsilon_{\mathbf{k}}$ 分别是 Bloch 波函数和能带能量。$\mathbf{k}$ 是电子的格点动量，$s=\uparrow,\downarrow$ 是电子在 $S^{z}$ 方向的自旋。这里我们略去能带指标，并忽略自旋轨道耦合。

对 $N_{e}$ 个电子，只有当哈密顿量是单粒子哈密顿量的可分离求和时，Schrödinger 方程才能约化为一组能带结构方程 (1.1)：

\begin{equation*}
\mathcal {H} ^ {0} = \sum_ {i = 1} ^ {N _ {e}} H ^ {0} [ \nabla_ {i}, \mathbf {x} _ {i} ].
\tag{1.2}
\end{equation*}

(1.2) 的本征态是由 Slater 行列式构造的 Fock 态（见附录 A）：

\begin{equation*}
\Psi_ {[ \mathbf {k}, s ]} ^ {\mathrm{Fock}} (\mathbf {x} _ {1}, \dots , \mathbf {x} _ {N _ {e}}) = \det _ {ij} \left[ \phi_ {\mathbf {k} _ {i}, s _ {i}} (\mathbf {x} _ {j}) \right]
\tag{1.3}
\end{equation*}

本征能量为

\begin{equation*}
E _ {[ \mathbf {k} ]} = \sum_ {i = 1} ^ {N _ {e}} \epsilon_ {\mathbf {k} _ {i}}.
\tag{1.4}
\end{equation*}

在基态中最低的 $N_{e}$ 个态被占据，其最高能量称为费米能（Fermi energy）：

\begin{equation*}
\max \epsilon_ {\mathbf {k}} = \epsilon_ {F}.
\tag{1.5}
\end{equation*}

式 (1.5) 在 $\mathbf{k}$ 空间中定义了费米面（Fermi surface）。

包含电子间相互作用的哈密顿量是

\begin{equation*}
\mathcal {H} = \mathcal {H} ^ {0} + \frac {1}{2} \sum_ {i \neq j} v ^ {\mathrm{el-el}} (\mathbf {x} _ {i}, \mathbf {x} _ {j}).
\tag{1.6}
\end{equation*}

$v^{\mathrm{el-el}}$ 破坏了 (1.2) 的可分离性，使 $\mathcal{H}$ 的对角化比 $\mathcal{H}^0$ 困难得多。

通过修改离子势，可以把大部分 Coulomb 相互作用效应纳入 $H$ 的单粒子部分：

\begin{equation*}
\mathcal {H} ^ {0} \rightarrow \tilde {\mathcal {H}} ^ {0} = \sum_ {i = 1} ^ {N _ {e}} \left( H ^ {0} [ \nabla_ {i}, \mathbf {x} _ {i} ] + V ^ {\mathrm{ion}} (\mathbf {x} _ {i}) + v ^ {\mathrm{eff}} [ \mathbf {x} _ {i}, \rho ] \right),
\tag{1.7}
\end{equation*}

其中 $v^{\mathrm{eff}}$ 是基态密度 $\rho(\mathbf{x})$ 的泛函。文献中对 $v^{\mathrm{eff}}[\rho]$ 有各种近似方案（见文献注释）。对具体材料的多数能带结构计算都把 $v^{\mathrm{eff}}$ 取为局域的，即只依赖 $\mathbf{x}$ 处的 $\rho$（或许还有其空间导数）。而 $\rho$ 本身又通过求解 (1.7) 的基态自洽地确定。

剩余相互作用（residual interactions）为

\begin{equation*}
\tilde {v} _ {ij} = v ^ {\mathrm{el-el}} (\mathbf {x} _ {i}, \mathbf {x} _ {j}) - \left[ v ^ {\mathrm{eff}} (\mathbf {x} _ {i}) + v ^ {\mathrm{eff}} (\mathbf {x} _ {j}) \right] / N _ {e}.
\tag{1.8}
\end{equation*}

变换 $v^{\mathrm{el-el}} \to \tilde{v}$ 代表屏蔽（screening）。实际上，屏蔽是一个动力学过程，涉及以等离子体频率为尺度的集体电荷涨落。若等离子体频率高于所关心的激发能量，$\tilde{v}$ 可以取为瞬时的。在金属中，被屏蔽的相互作用在 Thomas--Fermi 屏蔽长度 $\lambda^{\mathrm{TF}}$ 内指数衰减。

忽略 $\tilde{v}_{ij}$ 的单电子近似方案，在预言许多材料的体能量与结构方面取得了巨大成功。然而，对于表现出磁性或超导电性的体系，剩余相互作用是决定性的。

$\tilde{v}_{ij}$ 的效应可以用 Feynman 图和时序 Green 函数做微扰计算。对某些图类的重新求和（如随机相位近似中的做法）可用于描述无相互作用体系向有序基态的不稳定性。这些方法在许多教科书中都有详细论述，其中一部分列于本章文献注释。本书要强调的技术，是专门为补充微扰论、处理强电子--电子相互作用而设计的。

\section{场与相互作用}

在多电子 Hilbert 空间中，使用强制实现所有态反对称性的二次量子化算符是方便的（见附录 A）。场算符 $\psi_{s}^{\dagger}(\mathbf{x})$ 产生一个局域在 $\mathbf{x}$ 处、自旋态为 $s$ 的粒子：

\begin{equation*}
\langle \mathbf {x} ' s ' | \psi_ {s} ^ {\dagger} (\mathbf {x}) | 0 \rangle = \delta_ {ss'} \delta (\mathbf {x} - \mathbf {x} ').
\tag{1.9}
\end{equation*}

利用 (A.10)，可以用任何正交归一的单粒子基底 $\{\phi\}$ 展开场算符：

\begin{equation*}
\psi_ {s} ^ {\dagger} (\mathbf {x}) = \sum_ {i} \phi_ {i} ^ {*} (\mathbf {x}) c _ {is} ^ {\dagger},
\tag{1.10}
\end{equation*}

其中 $c^{\dagger}, c$ 是反对易的产生与湮灭算符。于是场算符满足

\begin{equation*}
\{ \psi_ {s} ^ {\dagger} (\mathbf {x}), \psi_ {s '} (\mathbf {x} ') \} = \delta (\mathbf {x} - \mathbf {x} ') \delta_ {ss'}.
\tag{1.11}
\end{equation*}

局域密度算符定义为

\begin{align*}
\hat {\rho} (\mathbf {x}) & = \sum_ {s} \psi_ {s} ^ {\dagger} (\mathbf {x}) \psi_ {s} (\mathbf {x})
\tag{1.12}\\
& = \sum_ {ii's} \phi_ {i} ^ {*} (\mathbf {x}) \phi_ {i '} (\mathbf {x}) c _ {is} ^ {\dagger} c _ {i's}.
\tag{1.13}
\end{align*}

$\hat{\rho}$ 度量在 $\mathbf{x}$ 处找到任一自旋粒子的概率密度。注意局域密度算符一般含有非对角项 $i \neq i'$。$\hat{\rho}(\mathbf{x})$ 是 Schrödinger 密度算符的二次量子化表示，后者在坐标空间中的期望值为

\begin{equation*}
\rho (\mathbf {x}) = \sum_ {i} \delta (\mathbf {x} - \mathbf {x} _ {i}).
\tag{1.14}
\end{equation*}

由 $|x_{1}, x_{2}, \ldots\rangle$ 规定的 Fock 态是 $\hat{\rho}(\mathbf{x})$ 的本征态，本征值为 $\rho(\mathbf{x})$。因此，用局域场算符与密度算符 (1.10) 和 (1.13) 来表示局域算符 $H$ 是自然的。像 (1.7) 与 (1.8) 那样把单粒子部分与剩余相互作用部分分开，我们写

\begin{equation*}
\mathcal {H} = \tilde {\mathcal {H}} ^ {0} + \tilde {\mathcal {V}} ^ {\mathrm{el-el}},
\tag{1.15}
\end{equation*}

其中

\begin{equation*}
\tilde {\mathcal {H}} ^ {0} = \sum_ {s} \int d ^ {3} x \, \psi_ {s} ^ {\dagger} (\mathbf {x}) \left[ - \frac {\hbar^ {2}}{2 m} \nabla^ {2} + V ^ {\mathrm{ion}} + v ^ {\mathrm{eff}} (\mathbf {x}) \right] \psi_ {s} (\mathbf {x}),
\tag{1.16}
\end{equation*}

以及

\begin{align*}
\tilde {\mathcal {V}} ^ {\mathrm{el-el}} & = \frac {1}{2} \int d ^ {3} x \, d ^ {3} y \, \tilde {v} (\mathbf {x}, \mathbf {y}) \left[ \hat {\rho} (\mathbf {x}) \hat {\rho} (\mathbf {y}) - \delta (\mathbf {x} - \mathbf {y}) \hat {\rho} (\mathbf {x}) \right]
\tag{1.17}\\
& = \frac {1}{2} \int d ^ {3} x \, d ^ {3} y \, \tilde {v} (\mathbf {x}, \mathbf {y}) \sum_ {ss'} \psi_ {s} ^ {\dagger} (\mathbf {x}) \psi_ {s '} ^ {\dagger} (\mathbf {y}) \psi_ {s '} (\mathbf {y}) \psi_ {s} (\mathbf {x}).
\end{align*}

在 (1.17) 的第二行中，场算符已正规排序，这消除了自相互作用项 $\tilde{v}(0)\,\delta(\mathbf{x}-\mathbf{y})\hat{\rho}$。

\begin{figure}[H]
\centering
\includegraphics[width=0.86\textwidth]{fig1_1a.jpg}\\[2pt]
\includegraphics[width=0.86\textwidth]{fig1_1b.jpg}
\caption{传导电子密度：(a)~$s$、$p$ 电子；(b)~$d$、$f$ 电子。实线是有效离子势。}
\label{fig:1.1}
\end{figure}

\section{金属中相互作用的量级}

必须强调，先验地判定一个给定体系中的剩余相互作用应当算“弱”还是“强”并不容易。恰当的问法是：它们对基态关联和低能激发的影响是否剧烈。微扰论可以作为估计相互作用效应的指南。粗略地说，金属中 Coulomb 相互作用效应的无量纲参数是

\begin{equation*}
g = \frac {e ^ {2}}{\kappa \epsilon_ {F} \left\langle r _ {ij} \right\rangle}
\tag{1.18}
\end{equation*}

其中 $e$ 是电子电荷，$\kappa$ 是静态介电常数，$\langle r_{ij}\rangle$ 是电子间平均距离，$\epsilon_{F}$ 是从导带底算起的费米能。贡献传导电子的离子在元素周期表中的位置越靠下，相对相互作用强度越大。离子势对 $s$、$p$ 传导电子的作用较弱，因此导带波函数是离域的，密度在整个晶体中相当均匀，如图~\ref{fig:1.1}(a) 所示。其结果是电子间相互作用相对较弱。

相反，过渡金属和混合价稀土化合物向导带贡献 $d$、$f$ 电子。在那里，电子大部分局域在围绕离子、半径 $\langle r_{ij}\rangle \ll a$ 的小区域内，$a$ 是晶格间距。同时，如图~\ref{fig:1.1}(b) 所示，大的原子间势垒压低了费米能 $\epsilon_{F}$。于是可以预期，过渡金属和混合价稀土化合物的 $g$ 值很大，应视为强相互作用电子体系。

这种分类只是粗略的经验法则。它无法预言某个具体化合物是顺磁金属，还是具有磁有序、电荷密度波或超导序。显然，无相互作用的简并基态对解除简并的微扰极为敏感。因此基态灵敏地依赖于能带结构的细节以及 $\tilde{H}^{0}$ 中的近简并。

\section{有效模型}

$\tilde{V}^{\mathrm{el-el}}$ 中的四费米子相互作用使 $H$ 成为真正的多体问题，因而极难对角化。数值方法只限于非常小的体系，因为 Hilbert 空间的维数随电子数目和单粒子基底的大小指数增长。理论家们于是求助于简化的模型（见第一部分扉页插图），其中只保留一个约减的单粒子态集合和相互作用矩阵元。

把 $H$ 替换为有效模型，若表述为重整化群变换\footnote{可参看 Shankar 的综述。}，则可以建立在更稳固的基础上。当人们关心低频、低温关联时，可以把 Green 函数投影到低能部分，从而消去高能态，得到一个只在低能子空间中作用的有效哈密顿量 $H^{\mathrm{eff}}$\footnote{例如：3.2 节中 t--J 模型就是作为大 $U$ Hubbard 模型的有效哈密顿量导出的。}。在路径积分表述中，可以把高频模式积掉，留下低频模式的有效 Lagrange 量\footnote{例如参看第 13 章中非线性 $\sigma$ 模型的重整化。}。

变换 $H \rightarrow H^{\mathrm{eff}}$ 称为重整化（renormalization）。在重整化下，某些相互作用（称为无关的，irrelevant）相对于其他相互作用被压低；增长或保持不变的分别被称为有关的（relevant）或边缘的（marginal）。这样，能带结构和相互作用的许多微观细节都从有效哈密顿量中消失，后者只包含最有关的相互作用。如果重整化得到一个单粒子态更少、相互作用更简化的模型，那么求低能关联就大有希望。有效模型的参数原则上可以通过求解重整化群流从微观哈密顿量确定。诚然，对真实材料的这类计算很少完成过，模型参数通常是拟合实验确定的。在缺乏可靠计算的情况下，某些相互作用无关还是有关，成了不同模型的支持者之间悬而未决的争论根源。

相互作用电子的一个最小模型是 Hubbard 模型，将在第 3 章描述。然而，即使是 Hubbard 模型也只在一维情形被严格求解。一个更简单、但仍包含多体相互作用的模型是量子 Heisenberg 模型，它将在第 3 章中作为 Hubbard 模型的一个特殊极限导出。虽然它只描述自旋自由度，但其基态和激发是高度关联的（即远非少数几个 Fock 态的组合）。Heisenberg 模型丰富的物理性质统称“量子磁性”（Quantum Magnetism），这是本书的主题之一。我们将用相当的篇幅讨论自旋相互作用及其对量子涨落和热涨落的影响。在更宽的视野里，量子 Heisenberg 模型将被用来演示基本的物理概念；作为一种教学工具，我们将在它上面检验各种对其他量子多粒子体系同样适用的数学技巧和近似方案。

\begin{exercises}

\item 利用附录 A 的 (A.24)，求自由玻色子和自由费米子的总粒子数 $N = -\partial F/\partial \mu$、熵 $S = -\partial F/\partial T$ 和比热 $C_{v} = T \,\partial S/\partial T$。

\end{exercises}

\begin{chapbib}
\item 关于固体电子结构的一般背景有大量优秀教科书，例如：
  N. W. Ashcroft and N. D. Mermin, \textit{Solid State Physics} (Holt, Rinehart, and Winston, 1976)；
  C. Kittel, \textit{Quantum Theory of Solids} (John Wiley and Sons, 1987)；
  J. Callaway, \textit{Quantum Theory of the Solid State} (Academic Press, 1992)；
  W. Jones and N. H. March, \textit{Theoretical Solid State Physics} (Dover, 1973)。
\item 凝聚态多体图微扰论的基本技术可参看，例如：
  A. Fetter and J. D. Walecka, \textit{Quantum Theory of Many Particle Systems} (McGraw-Hill, 1971)；
  G. Mahan, \textit{Many Particle Physics} (Plenum Press, 1981)。
\item P. Fulde, \textit{Electron Correlations in Molecules and Solids} (Springer-Verlag, 1991)。
  这本新著也值得推荐为当代理论与实验参考书。
\item 建议的进一步背景阅读：
  P. W. Anderson, \textit{Concepts in Solids} (Addison-Wesley, 1963)。
\item 关于相互作用电子重整化群的一篇出色新近教程：
  R. Shankar, Rev. Mod. Phys. \textbf{66}, 129 (1994)。
\item 关于相互作用电子的场论方法，两部著名的教科书是：
  E. Fradkin, \textit{Field Theories of Condensed Matter Systems} (Addison-Wesley, 1991)；
  A. M. Tsvelik, \textit{Quantum Field Theory in Condensed Matter Physics} (Cambridge University Press, 1995)。
\end{chapbib}
